\documentclass[11pt]{article}
\usepackage[margin=1in]{geometry}
\usepackage{amsmath,amssymb}
\usepackage{enumitem}
\usepackage{bookmark}
\hypersetup{
  pdftitle={PSET 5: Functions of several variables and optimization with several variables},
  hidelinks}

\title{PSET 5: Functions of several variables and optimization with
several variables}
\author{}
\date{\vspace{-2.5em}}


\newcounter{problem}
\newcommand{\problem}[1]{\stepcounter{problem}\section*{Problem \theproblem: #1}}

\begin{document}
\maketitle

% Shared assignment questions (header block).
% Include at the top of any problem set with:  % Shared assignment questions (header block).
% Include at the top of any problem set with:  % Shared assignment questions (header block).
% Include at the top of any problem set with:  \input{assignment-qs}
\section*{Assignment Qs}

\begin{itemize}
\item Name
\item How long did this problem set take you (round to nearest hour)? \quad
  0-1 hour \ 2-3 hours \ 4-5 hours \ 5+ hours
\item How difficult was this problem set? \quad
  very easy \ 1 \ 2 \ 3 \ 4 \ 5 \ very challenging
\end{itemize}

\section*{Assignment Qs}

\begin{itemize}
\item Name
\item How long did this problem set take you (round to nearest hour)? \quad
  0-1 hour \ 2-3 hours \ 4-5 hours \ 5+ hours
\item How difficult was this problem set? \quad
  very easy \ 1 \ 2 \ 3 \ 4 \ 5 \ very challenging
\end{itemize}

\section*{Assignment Qs}

\begin{itemize}
\item Name
\item How long did this problem set take you (round to nearest hour)? \quad
  0-1 hour \ 2-3 hours \ 4-5 hours \ 5+ hours
\item How difficult was this problem set? \quad
  very easy \ 1 \ 2 \ 3 \ 4 \ 5 \ very challenging
\end{itemize}


\problem{Evaluating multivariate functions}

Recall the spatial model of policy from class. Policy is $N$-dimensional,
$\mathbf{x} \in \mathbb{R}^{N}$, and legislator $i$'s utility
$U:\mathbb{R}^{N} \rightarrow \mathbb{R}^{1}$ is
\[
  U(\mathbf{x}) = -\sum_{j=1}^{N} (x_{j} - \mu_{j})^{2},
\]
where $\boldsymbol{\mu} = (\mu_1, \ldots, \mu_N)$ is the legislator's ideal
point. Suppose $N = 3$ and $\boldsymbol{\mu_1} = (0, 0, 0)$.

\begin{enumerate}[label=\alph*.]
\item (0.5 points) Evaluate $U(\mathbf{m})$ for the policy proposal
  $\mathbf{m_1} = (1, 2, -1)$, and for the proposal
  $\mathbf{m_2} = (0, 1, 0)$. Show your work.
\item (0.5 points) Which proposal does the legislator prefer, and how do you
  know from the values you just computed? 
\item (0.5 points) (in words only) How could you identify the legislator's preferred policy computationally, if you hadn't been told $\mu$ in advance. Be specific using terms and concepts from class. 
\end{enumerate}

\problem{First and second partial derivatives}

Find the requested partial derivatives of each
function.\footnote{Inspired by Grimmer HW6.3}

\begin{enumerate}[label=\alph*.]
\item (0.5 points) All first partial derivatives of $f(x,y) = \frac{1}{3}x^5 - 3x^3y^2 + 3xy^4$
\item (0.5 points) All first partial derivatives of $g(x,y) = xe^{4y}$
\item (0.5 points) All first partial derivatives of $k(x,y) = \dfrac{x+y}{x-y}$
\item (0.5 points) Returning to $f$ from part (a), find the second partial
  derivatives $\dfrac{\partial^2 f}{\partial x^2}$ and
  $\dfrac{\partial^2 f}{\partial y^2}$, and both cross partials
  $\dfrac{\partial^2 f}{\partial x \, \partial y}$ and
  $\dfrac{\partial^2 f}{\partial y \, \partial x}$. What do you notice about
  the two cross partials?
\end{enumerate}

\problem{Gradients, Hessians, and Jacobians}

\begin{enumerate}[label=\alph*.]
\item $g(x,y) = x^4 - 3x^2 y^3$
  \begin{enumerate}[label=\roman*.]
  \item (0.5 points) Find the gradient $\nabla g$ and evaluate it at
    $(x,y) = (1,1)$.
  \item (0.5 points) Find the Hessian $\mathbf{H}(g)$.
  \end{enumerate}
\item $f(x,y) = x^3y^4 - 3xy + 1$
  \begin{enumerate}[label=\roman*.]
  \item (0.5 points) Find the gradient $\nabla f$ and evaluate it at
    $(x,y) = (0,0)$.
  \item (0.5 points) Find the Hessian $\mathbf{H}(f)$.
  \end{enumerate}
\item (0.5 points) The gradient is the derivative of a \emph{scalar-valued}
  function; the Jacobian is the matrix of derivatives of a
  \emph{vector-valued} function. Consider
  $\mathbf{F}: \mathbb{R}^2 \rightarrow \mathbb{R}^2$ given by
  \[
    \mathbf{F}(x,y) = \left( x^2 y, \; x + y^3 \right).
  \]
  Find the Jacobian $\mathbf{J}(\mathbf{F})$. Then, in one sentence, explain
  why it would not be meaningful to ask for the Jacobian of $g$ in part (a)
  in the same sense.
\end{enumerate}

\problem{Find the critical points}

(1.5 points) Find the local minimum values, local maximum values, and saddle
point(s) of
\[
  f(x,y) = x^3 + y^3 - 3xy.
\]
Follow the process we discussed in class: calculate the gradient, set it
equal to zero and solve the resulting system of equations, calculate the
Hessian, and assess the Hessian at each critical value using the determinant
test. Be sure to show your work on each of these steps, and state clearly
what kind of point each critical value is.\footnote{Inspired by Grimmer HW7.4}

\problem{Definite and indefinite integrals}

Solve the following definite or indefinite integrals using the antiderivative
method.\footnote{Inspired by Gill 5.10, Grimmer HW4.1, Gill 5.13, and 5.14}

Reminder: the basic approach to compute the definite
integral of $f(x)$ from $a$ to $b$ is by using the formula $F(b) - F(a)$,
where $F(x)$ is the \textbf{antiderivative} of $f$.

\begin{enumerate}[label=\alph*.]
\item (0.5 points) $\int_{-1}^{0} (3x^2 - 1) \,dx$
\item (0.5 points) $\int_2^4 e^y \,dy$
\item (0.5 points) $\int_3^3 \sqrt{x^5 + 2} \,dx$
\item (0.5 points) $\int \left(x^2 - x^{-\frac{1}{2}}\right) \,dx$
\end{enumerate}

\problem{Substitution and integration by parts}

\begin{enumerate}[label=\alph*.]
\item (0.5 points) $\int 2x(x^2+1)^4 \,dx$
\item (0.5 points) $\int \dfrac{3x^2}{x^3+5} \,dx$
\item (0.5 points) $\int x e^{2x} \,dx$
\item (0.5 points) $\int \log(x) \,dx$ \quad (\emph{Hint}: you may take
  $dv = dx$.)

\item (0.5 points) Without solving, state whether you would use substitution or
  integration by parts for each, and why:
  \begin{enumerate}[label=\roman*.]
  \item $\int 5x^4 \left(x^5 - 2\right)^{3} \,dx$
  \item $\int x^2 \log(x) \,dx$
  \end{enumerate}
\end{enumerate}

\problem{Improper and multiple integrals}

For (a)--(c), determine whether the integral is convergent or divergent, and
evaluate those that are convergent.\footnote{Inspired by Grimmer HW4.3} Note
that these are improper for different reasons --- say which in each case.
For (d)--(e), evaluate the iterated integral.

\begin{enumerate}[label=\alph*.]
\item (0.5 points) $\int_1^{\infty} \left(\dfrac{1}{3x}\right)^2 \,dx$
\item (0.5 points) $\int_{-\infty}^0 x^3 \,dx$
\item (0.5 points) $\int_{0}^{1} x^{-\frac{1}{2}} \,dx$
\item (0.5 points) $\int_{0}^1 \int_{2}^{3} x^2y^3 \,dx \,dy$. Does the order
  in which you integrate matter here? Explain why or why not.
\item (0.5 points) $\int_{0}^1 \int_0^{\sqrt{1-x^2}} 2x^3y \,dy \,dx$. State
  which variable you must integrate first and why --- contrast your answer
  with part (d).
\end{enumerate}


\problem{Interpretation}

No new computation is required for this problem. Answer each part in two to
four sentences of prose.

\begin{enumerate}[label=\alph*.]
\item (1 point) \textbf{Reading a partial derivative.} Suppose we regress presidential
  approval in month $i$ on the employment rate and the price of gas, and
  obtain
  \[
    \text{Approval}_{i} = 0.8 - 0.5\,\text{Employ}_{i} - 0.25\,\text{Gas}_{i}.
  \]
  Write down $\dfrac{\partial\, \text{Approval}_{i}}{\partial\, \text{Employ}_{i}}$
  and interpret it substantively. Your interpretation should make clear what
  is happening to $\text{Gas}_{i}$ while employment changes, and why that
  matters for the claim you are making.

\item (1 point) \textbf{Gradient versus Hessian.} Explain why the condition
  $\nabla f(\mathbf{a}) = \mathbf{0}$ is \emph{necessary} but not
  \emph{sufficient} for $\mathbf{a}$ to be a local extremum. Use your answers
  to Problem 4 to illustrate: both $(0,0)$ and $(1,1)$ have zero gradient,
  but they are different kinds of point. What does the Hessian tell you at
  each that the gradient could not? Finally, explain what it means when the
  determinant test returns $AC - B^2 = 0$, and why we cannot classify the
  point in that case.

\item (0.5 points) \textbf{Integrals and probability.} A continuous probability density
  $f$ must satisfy $\int_{-\infty}^{\infty} f(x) \,dx = 1$. Explain what this
  requirement means, and what $\int_{a}^{b} f(x)\,dx$ represents for some
  interval $[a,b]$.

\end{enumerate}

% Shared AI and Resources statement.
% Include in any problem set with:  % Shared AI and Resources statement.
% Include in any problem set with:  % Shared AI and Resources statement.
% Include in any problem set with:  \input{ai-statement}
\section*{AI and Resources statement}

Please list (in detail) all resources you used for this assignment. If
you worked with people, list them here as well. It is not enough to say
that you used a resource for help, you need to be specific on the link
and \emph{how} it was helpful. W/R/T gen AI tools (including GPT,
Claude, etc.) you cannot use them to do work on your behalf---you
cannot put in any of the questions, etc. You can ask for help on logic
/ sample problems. If you do use GPT or other AI tools, you need to
provide a link to your chat transcript. Any suspected academic
integrity violations will be immediately reported to the Dean of
Students.

\section*{AI and Resources statement}

Please list (in detail) all resources you used for this assignment. If
you worked with people, list them here as well. It is not enough to say
that you used a resource for help, you need to be specific on the link
and \emph{how} it was helpful. W/R/T gen AI tools (including GPT,
Claude, etc.) you cannot use them to do work on your behalf---you
cannot put in any of the questions, etc. You can ask for help on logic
/ sample problems. If you do use GPT or other AI tools, you need to
provide a link to your chat transcript. Any suspected academic
integrity violations will be immediately reported to the Dean of
Students.

\section*{AI and Resources statement}

Please list (in detail) all resources you used for this assignment. If
you worked with people, list them here as well. It is not enough to say
that you used a resource for help, you need to be specific on the link
and \emph{how} it was helpful. W/R/T gen AI tools (including GPT,
Claude, etc.) you cannot use them to do work on your behalf---you
cannot put in any of the questions, etc. You can ask for help on logic
/ sample problems. If you do use GPT or other AI tools, you need to
provide a link to your chat transcript. Any suspected academic
integrity violations will be immediately reported to the Dean of
Students.


\end{document}
